\documentclass{report}
\usepackage{amsmath,amssymb}
\setlength{\parindent}{0mm}
\setlength{\parskip}{1em}
\begin{document}
\begin{center}
\rule{6in}{1pt} \
{\large C. T. Kelley \\
{\bf Sparse Interpolatory Models for Molecular Dynamics}}

NC State University \\ Dept of Mathematics \\ Box 8205 \\ Raleigh \\ NC 27695-8205
\\
{\tt tim\_kelley@ncsu.edu}\\
David Mokrauer\end{center}

We describe a method for using interpolatory models to accurately and
efficiently simulate molecular excitation and relaxation. We use
sparse interpolation for efficiency and local error estimation and
control for robustness and accuracy.

The objective of the project is to design an efficient
algorithm for simulation of light-induced molecular transformations.
The simulation seeks to follow the relaxation path of a molecule
after excitation by light. The simulator is a predictive tool to
see if light excitation and subsequent return to the unexcited
or ground state will produce a different configuration than
the initial one.

We simulate the results of the excitation, rather than the excitation
itself. The excitation will change the quantum state of a molecule. Our
objective is to design software that will let one explore the possible
changes in a molecule after a sequence of excitations.

The goals of the simulation are not only
to identify the end point, but to report the entire path in
an high-dimensional configuration space so that
one can look for nearby paths to interesting configurations and examine
the energy landscape near the path to see if low energy barriers make
jumping to a different path possible.


\end{document}
